%% LyX 2.3.6 created this file.  For more info, see http://www.lyx.org/.
%% Do not edit unless you really know what you are doing.
\documentclass[11pt,english]{article}
\usepackage[T1]{fontenc}
\usepackage[latin9]{inputenc}
\usepackage[a4paper]{geometry}
\geometry{verbose,tmargin=2cm,bmargin=2cm,lmargin=2cm,rmargin=2cm}
\usepackage{color}
\usepackage{babel}
\usepackage{amsmath}
\PassOptionsToPackage{normalem}{ulem}
\usepackage{ulem}
\usepackage[unicode=true]
 {hyperref}

\makeatletter

%%%%%%%%%%%%%%%%%%%%%%%%%%%%%% LyX specific LaTeX commands.
%% Because html converters don't know tabularnewline
\providecommand{\tabularnewline}{\\}

\makeatother

\begin{document}
\noindent \begin{flushleft}
\begin{tabular}{llr}
Mattias Villani  & \hspace{8.5cm} & Bayesian Learning \tabularnewline
Department of Statistics &  & \tabularnewline
Stockholm University  &  & \tabularnewline
\end{tabular}
\par\end{flushleft}

\vspace{0.3cm}

\noindent \begin{flushleft}
\begin{tabular}{ll}
\multicolumn{2}{l}{\textbf{\large{}{}Computer Lab 2}}\tabularnewline
\multicolumn{2}{l}{\rule{0.975\columnwidth}{1pt}}\tabularnewline
\multicolumn{2}{l}{The labs should be done in R, since only R is available at the computer
exam. }\tabularnewline
Your lab solutions will be available for you at the exam.  & \tabularnewline
\multicolumn{2}{l}{You work and submit your labs in pairs, but both of you should contribute
equally}\tabularnewline
and understand all parts of your solutions. & \tabularnewline
\multicolumn{2}{l}{\texttt{\textbf{\textcolor{blue}{It is not allowed to share exact
solutions}}} with other student pairs.}\tabularnewline
Submit your solutions via Athena. & \tabularnewline
The datasets are available in the package \texttt{SUdatasets}. See
\textcolor{blue}{\uline{\href{https://github.com/mattiasvillani/BayesLearnCourse\#computer-labs}{the course page}}}
for instructions. & \tabularnewline
\multicolumn{2}{l}{\rule{0.975\columnwidth}{1pt}}\tabularnewline
\end{tabular}
\par\end{flushleft}

{\footnotesize{}\vspace{0.5cm}
}{\footnotesize\par}
\begin{enumerate}
\item \textbf{Linear and polynomial regression}\\
The dataset $\texttt{tempLinkoping}$ in the package \texttt{SUdatasets}
contains daily temperatures (in Celcius degrees) at Malmsl�tt, Link�ping
over the course of the year 2016 (366 days since 2016 was a leap year).
The response variable is $temp$ and the covariate is
\[
time=\frac{\text{the number of days since beginning of year}}{366}.
\]
The task is to perform a Bayesian analysis of a quadratic regression
\[
temp=\beta_{0}+\beta_{1}\cdot time+\beta_{2}\cdot time^{2}+\varepsilon\mbox{,\text{ }}\varepsilon\overset{iid}{\sim}N(0,\sigma^{2}).
\]

\begin{enumerate}
\item \textit{Determining the prior distribution of the model parameters.}
Your task is to set the hyperparameters $\mu_{0}$, $\Omega_{0}$,
$\nu_{0}$ and $\sigma_{0}^{2}$ in the conjugate prior to sensible
values. Start with $\mu_{0}=(-10,100,-100)^{T}$, $\Omega_{0}=0.01\cdot I_{3}$,
$\nu_{0}=4$ and $\sigma_{0}^{2}=1$. Check if this prior agrees with
your prior opinions by simulating draws from the joint prior of all
parameters and for every draw compute the regression curve. This gives
a collection of regression curves, one for each draw from the prior.
Do the collection of curves look reasonable? If not, change the prior
hyperparameters until the results agree with your prior beliefs about
the regression curve.\\
{[}Hint: the R package $\texttt{mvtnorm}$ will be handy. And use
your $Inv$-$\chi^{2}$ simulator from Lab 1.{]}
\item Write a program that \textit{simulates from the joint posterior distribution}
of $\beta_{0}$, $\beta_{1}$,$\beta_{2}$ and $\sigma^{2}$. Plot
the marginal posteriors for each parameter as a histogram. Also produce
another figure with a scatter plot of the temperature data and overlay
a curve for the posterior median of the regression function $f(time)=\beta_{0}+\beta_{1}\cdot time+\beta_{2}\cdot time^{2}$,
computed for every value of $time$. Also overlay curves for the lower
2.5\% and upper 97.5\% posterior credible interval for $f(time)$.
That is, compute the 95\% equal tail posterior probability intervals
for every value of $time$ and then connect the lower and upper limits
of the interval by curves. Does the interval bands contain most of
the data points? Should they?
\item It is of interest to locate the $time$ with the highest expected
temperature (that is, the $time$ where $f(time)$ is maximal). Let's
call this value $\tilde{x}$. Use the simulations in b) to simulate
from the \textit{posterior distribution of} $\tilde{x}$. {[}Hint:
the regression curve is a quadratic. You can find a simple formula
for $\tilde{x}$ given $\beta_{0},\beta_{1}$ and $\beta_{2}$.{]}
\item Say now that you want to \textit{estimate a polynomial model of order
$7$}, but you suspect that higher order terms may not be needed,
and you worry about over-fitting. Suggest a suitable prior that mitigates
this potential problem. You do not need to compute the posterior,
just write down your prior. {[}Hint: the task is to specify $\mu_{0}$
and $\Omega_{0}$ in a smart way.{]}
\end{enumerate}
{\footnotesize{}\vspace{0.5cm}
}{\footnotesize\par}
\item \textbf{Posterior approximation for classification with logistic regression}\\
The dataset \texttt{womenwork} in the package \texttt{SUdatasets}
contains $n=200$ observations on the following nine variables:\\
 %
\begin{tabular}{lllllll}
 &  &  &  &  &  & \tabularnewline
\cline{4-7} \cline{5-7} \cline{6-7} \cline{7-7} 
 &  &  & Variable & Data type & Meaning & Role\tabularnewline
\cline{4-7} \cline{5-7} \cline{6-7} \cline{7-7} 
 &  &  & \texttt{work}  & Binary & Whether or not the woman works & Response\tabularnewline
 &  &  & \texttt{constant } & 1 & Constant to the intercept & Feature\tabularnewline
 &  &  & \texttt{husbandInc} & Numeric & Husband's income & Feature\tabularnewline
 &  &  & \texttt{educYears} & Counts & Years of education & Feature\tabularnewline
 &  &  & \texttt{expYears} & Counts & Years of experience & Feature\tabularnewline
 &  &  & \texttt{expYears2} & Numeric & (Years of experience/10)$^{2}$ & Feature\tabularnewline
 &  &  & \texttt{age } & Counts & Age & Feature\tabularnewline
 &  &  & \texttt{nSmallChild} & Counts & Number of child $\leq6$ years & Feature\tabularnewline
 &  &  & \texttt{nBigChild} & Counts & Number of child $>6$ years & Feature\tabularnewline
\cline{4-7} \cline{5-7} \cline{6-7} \cline{7-7} 
\end{tabular}

\begin{enumerate}
\item Consider the logistic regression
\[
\mathrm{Pr}(y=1\vert\mathbf{x})=\frac{\exp\left(\mathbf{x}^{T}\beta\right)}{1+\exp\left(\mathbf{x}^{T}\beta\right)},
\]
where $y$ is the binary variable with $y=1$ if the woman works and
$y=0$ if she does not. $\mathbf{x}$ is a $8$-dimensional vector
containing the eight features (including a one for the constant term
that models the intercept). Fit the logistic regression using maximum
likelihood estimation by the command: \texttt{glmModel <- glm(work
\textasciitilde{} 0 + ., data = womenwork, family = binomial)}. Note
how I added a zero in the model formula so that R doesn't add an extra
intercept (we already have an intercept term from the \texttt{constant}
feature). Note also that a dot (.) in the model formula means to add
all other variables in the dataset as features. \texttt{family = binomial
}tells R that we want to fit a logistic regression.
\item Now the fun begins. Our goal is to approximate the posterior distribution
of the 8-dim parameter vector $\beta$ with a multivariate normal
distribution
\[
\beta\vert\mathbf{y},\mathbf{X}\sim N\left(\tilde{\beta},J_{\mathbf{y}}^{-1}(\tilde{\beta})\right),
\]
where $\tilde{\beta}$ is the posterior mode and $J(\tilde{\beta})=-\frac{\partial^{2}\ln p(\beta\vert\mathbf{y})}{\partial\beta\partial\beta^{T}}\vert_{\beta=\tilde{\beta}}$
is the observed Hessian evaluated at the posterior mode. Note that
$\frac{\partial^{2}\ln p(\beta\vert\mathbf{y})}{\partial\beta\partial\beta^{T}}$
is an $8\times8$ matrix with second derivatives on the diagonal and
cross-derivatives $\frac{\partial^{2}\ln p(\beta\vert\mathbf{y})}{\partial\beta_{i}\partial\beta_{j}}$
on the off-diagonal. Now, both $\tilde{\beta}$ and $J(\tilde{\beta})$
are computed by the \texttt{optim} function in R. You can use the
code in this \textcolor{blue}{\uline{\href{https://github.com/mattiasvillani/BayesLearnCourse/raw/master/Notebooks/R/SpamOptim.Rmd}{R notebook}}}
as a guide, but don't copy my code, write your own. Use the prior
$\beta\sim\mathcal{N}(0,\tau^{2}I)$, with $\tau=10$.\\
Your report should include your code as well as numerical values for
$\tilde{\beta}$ and $J_{\mathbf{y}}^{-1}(\tilde{\beta})$ for the
\texttt{womenwork} data. Compute an approximate 95\% credible interval
for the coefficient on \texttt{nSmallChild}. Would you say that this
feature is an important determinant of the probability that a women
works?
\item Write a function that simulates from the predictive distribution of
the response variable in a logistic regression. Use your normal approximation
from 2(b). Use that function to simulate and plot the predictive distribution
for the \texttt{work} variable for a 40 year old woman, with two children
(3 and 9 years old), 8 years of education, 10 years of experience,
and a husband with an income of 10.\\
{[}Hint: the R package $\texttt{mvtnorm}$ will again be handy. And
remember my discussion on how Bayesian prediction can be done by simulation.{]}
\end{enumerate}
\end{enumerate}

\end{document}
